\BeleskeID{09_2-s002}
% element: 09_2-s002:e04
Ugrađeni kanal postoji i pri nultom naponu između gejta i sorsa. Prag opterećenja je negativan, pa uslov za provođenje ostaje ispunjen. Zakočeni pogonski tranzistor ne odvodi struju, a neopterećeni izlaz nema drugu putanju struje. Zato je i struja opterećenja jednaka nuli. U omskoj oblasti to se ostvaruje pri $V_{DS,L}=0$, odnosno izlaz dostiže napajanje.

Pretpostavka da opterećenje radi u zasićenju dala bi
\[I_{Dn,L}=\frac{k_{n,L}}2(V_{GS,L}-V_{Tn,L})^2=\frac{k_{n,L}}2(-V_{Tn,L})^2>0.\]
Takva struja nije saglasna sa zakočenim pogonskim tranzistorom i neopterećenim izlazom. Nulta struja stoga bira radnu tačku u omskoj oblasti, nezavisno od veličine prekoračenja praga.
